%! Author = sriadityavedantam
%! Date = 3/29/26

\section{Axioms}

Axioms are trivial definitions used in everyday life, or even mathematics, that we take for granted.
They are definitions that are inarguable and are the core of mathematics today.I never quite understood the hierarchy of math statements, but this is a way to look at it.
Axioms are a specific type of definition that is just taken as a fact or true.
Definitions are similar to axioms in that they build the premise of future statements, and these may or may not include proofs to explain why they are true.
Lemmas are true statements that are not important in the long run, but are useful for understanding future statements, and are generally associated with proofs.
Propositions are important statements that must be associated with proofs and are vital building blocks in research.
Theorems are big conclusions that wrap each concept mentioned in a paper into one central idea and are even more important than propositions, and these also require proofs to be stated alongside the statement.Now the following axioms or properties are what we accept without another thought, but they are important to mention to understand future content when they are brought up again.

\begin{definition}[Additive Properties]
    \begin{enumerate}
        \item Addition is well-defined.
        Given $a,b \in \mathbb{Z}$, $a + b$ is a uniquely defined integer.

        \item Substitution Law.
        Since addition is well-defined, if $a = b$ and $c = d$, then $a + c = b + d$.

        \item Commutative Law.
        For all $a,b \in \mathbb{Z}$, $a + b = b + a$.

        \item Associative Law.
        For all $a,b,c \in \mathbb{Z}$, $(a + b) + c = a + (b + c)$.

        \item There exists a zero element $0 \in \mathbb{Z}$, called the additive identity, satisfying $0 + a = a$ for any $a \in \mathbb{Z}$.

        \item For all $a \in \mathbb{Z}$, there exists a unique additive inverse $-a \in \mathbb{Z}$ satisfying $a + (-a) = 0$.
    \end{enumerate}
\end{definition}

\begin{definition}[Multiplicative Properties]
    \begin{enumerate}
        \item Multiplication is well-defined.
        Given $a,b \in \mathbb{Z}$, $a \cdot b$ is a uniquely defined integer.

        \item Substitution Law.
        If $a = b$ and $c = d$, then $ac = bd$.

        \item $\mathbb{Z}$ is closed under multiplication.
        For all $a,b \in \mathbb{Z}$, $a \cdot b \in \mathbb{Z}$.

        \item Commutative Law.
        For all $a,b \in \mathbb{Z}$, $ab = ba$.

        \item Associative Law.
        For all $a,b,c \in \mathbb{Z}$, $(ab)c = a(bc)$.

        \item There exists a multiplicative identity $1 \in \mathbb{Z}$ satisfying $1 \cdot a = a$ for all $a \in \mathbb{Z}$.
    \end{enumerate}
\end{definition}

\begin{definition}[Distributive Property]
    For all $a,b,c \in \mathbb{Z}$,
    \[
        a(b + c) = ab + ac.
    \]
\end{definition}

\begin{definition}[Trichotomy Principle]
    $\mathbb{Z}$ can be split into three distinct sets:
    \[
        \mathbb{Z} = -\mathbb{N} \cup \{0\} \cup \mathbb{N}.
    \]
\end{definition}

\begin{definition}[Positivity Axiom]
    The sum or product of positive integers is positive.
\end{definition}

\begin{definition}[Discrete Property]
    We have learned these already, namely the Well-Ordering Principle of $\mathbb{N}$ and the Principle of Induction.
\end{definition}